\clearpage\fancyhead[L]{\small PALLAS.}\hypertarget{native-373-lower}{}\label{native-373-lower}
\begin{center}\textbf{PALLAS.}\par\vspace{4pt}\hrule\vspace{4pt}\fontsize{10}{12}\selectfont Monatliche Correspondenz. Band XXVIII. Seite 197\;..\;198. August 1813.\par\vspace{4pt}\hrule\vspace{4pt}\end{center}

\begin{flushright}
Göttingen, 31. August 1813.
\end{flushright}

\par\indent --- --- --- Ich sende Ihnen hier meine diesjährigen Beobachtungen der \emph{Pallas}:

\begin{center}\fontsize{10}{12}\selectfont\setlength{\tabcolsep}{5pt}\renewcommand{\arraystretch}{1.15}\begin{tabular}{l|c|c|c}
1813 & \shortstack{Mittlere Zeit\\in Göttingen} & Rektasz. $\PallasSymbol$ & Dekl.\\\hline
Juni 28 & $11^{\mathrm u}\;37^{\mathrm m}\;2^{\mathrm s}$ & $327^\circ\;2'\;13''3$ & $+14^\circ\;5'\;43''4$\\
Aug. 16 & $10\;36\;23$ & $319\;50\;54.7$ & $.\;.\;.\;.\;.$\\
19 & $9\;31\;0$ & $319\;17\;10.6$ & $10\;0\;8.4$\\
26 & $9\;23\;7$ & $317\;59\;52.0$ & $8\;44\;35.3$\\
\end{tabular}\end{center}

\Needspace{12\baselineskip}\par\indent Der Vergleich von Herrn \textsc{Nicolai} mit den neuesten Elementen ergibt folgende Ergebnisse:
\begin{center}\fontsize{10}{12}\selectfont\setlength{\tabcolsep}{7pt}\renewcommand{\arraystretch}{1.15}\begin{tabular}{l|c|c}
\SourceLeaf{374} & \multicolumn{2}{c}{Unterschied}\\
1813 & in Rektasz. & in Dekl.\\\hline
Juni 28 & $+23''6$ & $+22''4$\\
Aug. 16 & $+35.6$ & $.\;.\;.\;.$\\
19 & $+28.5$ & $+4.5$\\
26 & $+22.4$ & $-4.3$\\
\end{tabular}\end{center}

\par\indent Ich bemerke dazu, dass die Rektaszension vom 16. August nur auf einem einzigen Vergleich beruht. Die Bestimmungen vom 19. und 26. August stützen sich dagegen auf zahlreiche Vergleiche, die ausgezeichnet übereinstimmen.


